\section{Problem, model, and representation classes}
\label{sec:problem}

\subsection{Three meanings of a finite dynamical description}

The word \emph{finite} hides substantial differences.  We use the following
classification.

\begin{definition}[Scalar, population, and operator state]
A \emph{scalar ODE} has state in $\R^q$ for fixed $q$.  A
\emph{population ODE} has a fixed finite list of scalar or vector fields on
fixed probability spaces; an $n$-particle discretization stores one value per
particle and field.  An \emph{operator ODE} has a fixed finite list of
operators between population spaces.  Such an operator can retain arbitrary
row--column correlations and is generally not specified by finitely many
population fields or scalars.
\end{definition}

Counting only the number of named objects would call all three descriptions
finite, but their information content is different.  A general limiting
operator is the continuum analogue of a dense matrix.  Likewise, one
unrestricted field can encode an arbitrary amount of information in its
spatial profile.  Any compression statement must therefore name both the
number and the type of retained objects.

Our closure is a finite collection of neuron-population fields indexed by
sample, layer, and history mode, together with fixed initialized operators.
The learned moving part is population-valued; the total representation is a
population/operator hybrid.  At finite width these fields are simply neuron
vectors.

Exact finite scalar closure is already obstructed in a broad natural class.
The following result is included to calibrate the target, not to claim an
impossibility theorem for every conceivable encoder.

\begin{proposition}[Restricted exact scalar nonclosure]
\label{prop:scalar-no-go}
Consider a state-universal deep linear feature-learning system with several
trainable factors.  Suppose an encoder consists of a fixed finite list of
complete tensor contractions whose connected graph degree is bounded
independently of width, its dynamics are a width-independent polynomial ODE,
and the output is a polynomial readout.  No such encoder exactly realizes the
output dynamics on a nonempty open set for all sufficiently large widths.
The same conclusion holds for analytic finite-jet closures near the zero
state.  It does not rule out arbitrary encoders or a closure restricted to one
initialized trajectory.
\end{proposition}

The proof, given in Appendix~\ref{app:scalar-no-go}, repeatedly differentiates
the output.  It produces connected contractions of unbounded graph degree,
whereas polynomial operations on a fixed bounded-degree alphabet can only
form disjoint unions of its existing connected components.  Stable graph
independence at large width makes the contradiction algebraic.  This result
motivates approximate population memory; it does not prove that every exact
scalar compression is impossible.

\subsection{Deep nonlinear feature-learning gradient flow}

Fix input dimension $d$, sample count $m$, width $n$, and $L\ge2$ hidden
layers.  The data are $(x_a,y_a)_{a=1}^m$ and $u_a=x_a/\sqrt d$.  Hidden
layers have common width only to simplify notation.  With coordinatewise
activations $\phi_\ell$,
\begin{align}
z_{1,a}&=W_1u_a, & h_{1,a}&=\phi_1(z_{1,a}),\nonumber\\
z_{\ell,a}&=W_\ell h_{\ell-1,a}, &
h_{\ell,a}&=\phi_\ell(z_{\ell,a}),\qquad 2\le\ell\le L,\label{eq:model-forward}\\
f_a&=\frac1n w^\top h_{L,a}, & r_a&=f_a-y_a .\nonumber
\end{align}
Here $W_1\in\R^{n\times d}$, $W_\ell\in\R^{n\times n}$ for
$2\le\ell\le L$, and $w\in\R^n$.  We use unhalved mean squared loss
\begin{equation}
 \cL=\rho^2,\qquad
 \rho=\left(\frac1m\sum_{a=1}^m r_a^2\right)^{1/2}.
 \label{eq:loss}
\end{equation}
Backpropagated response fields omit the loss derivative:
\begin{align}
\delta_{L,a}&=w\odot\phi_L'(z_{L,a}),\nonumber\\
\delta_{\ell,a}&=\phi_\ell'(z_{\ell,a})\odot
 W_{\ell+1}^\top\delta_{\ell+1,a},\qquad \ell<L .
\label{eq:backward}
\end{align}

The block mobilities $(n,1,\ldots,1,n)$ are the canonical
feature-learning scaling used throughout the paper.  The resulting gradient
flow is
\begin{align}
\dot W_1&=-\frac2m\sum_a r_a\delta_{1,a}u_a^\top,\nonumber\\
\dot W_\ell&=-\frac{2}{nm}\sum_a
 r_a\delta_{\ell,a}h_{\ell-1,a}^\top,
 \qquad 2\le\ell\le L,\label{eq:dense-gf}\\
\dot w&=-\frac2m\sum_a r_a h_{L,a}.\nonumber
\end{align}
All layers move at order one under the intended wide scaling, so the model is
not reduced to a frozen neural tangent kernel \citep{jacot2018ntk,yang2021tensor}.

Let $\theta=(W_1,\ldots,W_L,w)$ and denote the right-hand side of
\eqref{eq:dense-gf} by $F(\theta)$.  Integrating an internal-layer equation
gives \eqref{eq:intro-history}.  Define, when $\rho>0$,
\begin{equation}
 b_{\ell,a}=\frac{r_a\delta_{\ell,a}}{\rho}.
 \label{eq:b-definition}
\end{equation}
Then the learned increment is the integral of $b_{\ell,a}h_{\ell-1,a}^\top$
against activity measure $\rho(t)\od t$.  The closure will compress exactly
these paired histories.

\subsection{Why the population regime remains the target}

Under Gaussian or sub-Gaussian initialization, fixed depth and data, and
globally Lipschitz activation derivatives, established multilayer mean-field
theorems construct nonlinear population flows on a common local interval and
identify finite-network observables as width grows
\citep{nguyen2023rigorous}.  Maximal-update tensor programs likewise identify
feature-learning limits for discrete training \citep{yang2021tensor}.  These
results justify studying a deterministic population target rather than only
one finite random network.

The following specialized background statement records the population result
used to motivate this paper's regime.

\begin{proposition}[Local population-flow identification]
\label{prop:population-background}
Fix $d,m,L$, bounded data, and activations with bounded first derivative and
globally Lipschitz first derivative.  Initialize the input weights and readout
independently with fixed sub-Gaussian laws and each internal matrix with
independent entries of variance $\sigma_\ell^2/n$.  Under the canonical
feature-learning mobilities, there is a deterministic $T_*>0$ and a unique
autonomous strong population field/operator flow on $[0,T_*]$.  For every
deterministic gradient-descent mesh tending to zero as width grows,
predictions, loss, and fixed kernel entries converge in probability uniformly
in time; separately in each layer, empirical laws of the joint $m$-input
hidden paths converge in quadratic Wasserstein distance.  A diagonal
finite-width Euler argument gives the same population limit for finite-network
gradient flow.
\end{proposition}

The proof constructs initialized forward and adjoint operators jointly,
closes the local population integral equation by Picard iteration, and obtains
response and tail bounds uniform in the Euler mesh.  A reference-oracle network
uses the deterministic population residuals and contractions; concentration
controls that oracle, and a discrete stability inequality transports the
error to the empirical network.  Truncation and uniform integrability remove
the sub-Gaussian tails.  Finally, for each width one chooses an Euler mesh fine
enough to approximate finite gradient flow and diagonalizes with the width
limit.  The result is local in population time; it does not supply a
global-in-time population law or the width-uniform response-memory estimate
sought here.

The approximation theorem below deliberately proves something logically
different: for each realized finite network it controls the response closure
against that network's dense gradient flow, for every prescribed finite
horizon.  No random initialization is required.  This pathwise statement is
useful for implementation and is a prerequisite for, but not a substitute
for, a width-uniform population theorem.  Section~\ref{sec:population-bridge}
states the exact diagonal consequence currently available.
